% chapters/appendix-ch11-sql.tex
\chapter{SQL Appendix: Chapter 11 (State and Time)}

\section{DDL}
\begin{verbatim}
CREATE TABLE copy (
  copy_id INT PRIMARY KEY AUTO_INCREMENT,
  status VARCHAR(20) NOT NULL
) ENGINE=InnoDB;

CREATE TABLE loan (
  loan_id INT PRIMARY KEY AUTO_INCREMENT,
  copy_id INT NOT NULL,
  checkout_date DATE NOT NULL,
  return_date DATE NULL,
  CONSTRAINT fk_loan_copy FOREIGN KEY (copy_id) REFERENCES copy(copy_id)
) ENGINE=InnoDB;
\end{verbatim}

\section{Sample data}
\begin{verbatim}
INSERT INTO copy (status) VALUES ('available'), ('on_loan');
INSERT INTO loan (copy_id, checkout_date, return_date)
VALUES (2, '2026-02-01', NULL);
\end{verbatim}

\section{Example queries}
\begin{verbatim}
-- Q1: Derive active state from events
SELECT c.copy_id,
       CASE WHEN EXISTS (
         SELECT 1 FROM loan l
         WHERE l.copy_id = c.copy_id AND l.return_date IS NULL
       ) THEN 'on_loan' ELSE 'available' END AS derived_status
FROM copy c;

-- Q2: Detect mismatch if status is stored
SELECT c.copy_id, c.status
FROM copy c
WHERE c.status = 'available'
  AND EXISTS (SELECT 1 FROM loan l WHERE l.copy_id = c.copy_id AND l.return_date IS NULL);
\end{verbatim}
